\begin{afterword}
\summaryhead

The post asks about the decision theory of things one should not do even when they seem right. It starts with an AI whose goals are unfinished. If such an AI ever decides that deceiving its programmers is right, that is probably a sign its goals have gone wrong, so it should shut down automatically instead of recalculating. The author does not yet know the mathematics.

The human version is an ethical injunction, such as never murdering an innocent person who has helped you, because a mistake is far likelier than a real exception. The sinking of the ferry Hydro in 1944, which the author says ended the Nazi atomic program, is a case the author cannot condemn. For self-deception the author knows of no exception: self-deceptions are black swan bets, whose rare blowups cancel their benefits.

Injunctions rest on two grounds, corrupted hardware and black swans. To the objection that one could include these in a fresh calculation, the author replies that this is the cleverness to beware of, that meta-reasoning can be corrupted too, and that ethics have protected the author before, while granting that the answer is not automatically settled. The post then dismisses people who argue for relaxing their ethics, and declines to endorse absolute injunctions.

\responsehead

The post turns the view of the previous post in the book into a practical rule, and states it with care. It raises the standard objection to such rules itself: if you know why a rule exists, why not include those reasons in a new calculation? It gives three replies and then concedes that the answer to ``Am I still dumber than my ethics?'' is not automatically yes. It declines to endorse absolute injunctions. Its first example is a case the author cannot condemn, which fits that refusal. The idea has a long history in ethics, from Mill's remark that self-deception enters through the latitude every creed allows, to the two-level views noted under ``Ends Don't Justify Means (Among Humans)''; that history supports the post.

The historical example is accurate in the details I could check, and ends on an overstatement. The sinking did not end the German program, whose reactor experiments went on into 1945, and the ferry's cargo held too little heavy water to supply a reactor. The author's hedge that Germany ``very likely'' would not have got the bomb matches the historians' view.

The self-deception section argues from unknowable risks, which is reasonable, and adds one unhedged claim: that all the happiness belief in an afterlife has produced has been ``more than cancelled'' by the failure to adopt cryonics. Neither side is measured, and the claim depends on cryonics working, which has not been shown.

The last third changes method. It describes the people who argue for relaxing their ethics, from the author's experience: they ``seem to have no idea of the risks'', they move ``Always'' toward lenience, and their positions are answered with ``Ha'' and ``Double ha''. The argument is complete before this section, which adds characterization, not reasons.

In short: a careful case for keeping rules against one's own calculation, which raises the main objection itself and declines to claim absolute rules, followed by a closing third that dismisses its opponents instead of answering them.
\end{afterword}
